Hoje nos debruçaremos sobre uma questão elegante que conecta os conceitos de autovalores e a
estrutura de uma matriz. A proposição a ser demonstrada é um pilar na compreensão dos
operadores nilpotentes, e sua prova revela a beleza e a força do Teorema de Cayley-Hamilton.

\subsection*{Enunciado do Problema:}

Mostre que se uma matriz $A \in M_n(\mathbb{C})$ (uma matriz de ordem $n \times n$ com entradas complexas) tem
somente o escalar $0$ como autovalor, então $A$ é nilpotente. Isto é, existe um inteiro positivo $k$ tal que
$A^k = 0$, onde $0$ é a matriz nula de ordem $n \times n$.

\subsection*{Demonstração Canônica}

Para provar esta afirmação com o rigor que a nossa disciplina exige, seguiremos um caminho
axiomático, fundamentado em um dos teoremas mais profundos da Álgebra Linear.

\subsubsection*{Passo 1: O Polinômio Característico}

A pedra angular para o estudo de autovalores é o polinômio característico. Por definição, o polinômio
característico da matriz $A$, denotado por $p_A(t)$, é dado por:

\[
p_A(t) = \det(A - tI_n)
\]

onde $I_n$ é a matriz identidade de ordem $n$.

Sabemos, por um teorema fundamental, que as raízes do polinômio característico são precisamente
os autovalores da matriz $A$.

\subsubsection*{Passo 2: Aplicando a Hipótese do Problema}

A hipótese que nos foi dada é que o \textbf{único} autovalor de $A$ é $\lambda = 0$.

Se $0$ é a única raiz do polinômio $p_A(t)$ sobre o corpo dos números complexos $\mathbb{C}$, então, pelo
Teorema Fundamental da Álgebra, este polinômio deve ser fatorado de uma forma muito específica.
Um polinômio de grau $n$ que possui apenas a raiz $t=0$ com multiplicidade algébrica $n$ deve ter a
forma:

\[
p_A(t) = c \cdot (t-0)^n = c \cdot t^n
\]

onde $c$ é uma constante.

Precisamos determinar o valor desta constante $c$. Vamos expandir o determinante que define $p_A(t)$:

\[
p_A(t) = \det(A - tI_n) =
\begin{vmatrix}
a_{11} - t & a_{12} & \cdots & a_{1n} \\
a_{21} & a_{22} - t & \cdots & a_{2n} \\
\vdots & \vdots & \ddots & \vdots \\
a_{n1} & a_{n2} & \cdots & a_{nn} - t
\end{vmatrix}
\]

O termo de maior grau em $t$ neste polinômio é obtido pelo produto dos elementos da diagonal
principal:

\[
(a_{11} - t)(a_{22} - t)\cdots(a_{nn} - t) = (-1)^n t^n + \text{termos de grau inferior}
\]

Ao comparar o termo líder, $(-1)^n t^n$, com a forma $c \cdot t^n$, concluímos que a constante $c$ deve ser
$(-1)^n$.

Portanto, o polinômio característico da matriz $A$ é, inequivocamente:

\[
p_A(t) = (-1)^n t^n
\]

\subsubsection*{Passo 3: A Invocação do Teorema de Cayley-Hamilton}

O Teorema de Cayley-Hamilton é uma das joias da nossa área. Ele afirma que toda matriz quadrada
anula seu próprio polinômio característico. Formalmente:

\textbf{Teorema (Cayley-Hamilton):} Se $p_A(t)$ é o polinômio característico de uma matriz $A$, então
$p_A(A) = 0$.

Isso significa que se substituirmos a variável $t$ pela matriz $A$ na expressão do polinômio
característico, o resultado será a matriz nula.

\subsubsection*{Passo 4: A Conclusão da Prova}

Vamos agora aplicar este poderoso teorema ao nosso resultado do Passo 2. Temos:

\[
p_A(t) = (-1)^n t^n
\]

Segundo Cayley-Hamilton, devemos ter $p_A(A) = 0$. Substituindo $t$ por $A$:

\[
p_A(A) = (-1)^n A^n
\]

Logo, o teorema nos garante que:

\[
(-1)^n A^n = 0
\]

Como $(-1)^n$ é um escalar não nulo (pode ser $1$ ou $-1$), podemos multiplicar ambos os lados da
equação pelo seu inverso, que é ele mesmo, $(-1)^n$.

\[
(-1)^n \cdot [(-1)^n A^n] = (-1)^n \cdot 0
\]

\[
((-1)^{2n}) A^n = 0
\]

\[
1 \cdot A^n = 0
\]

\[
A^n = 0
\]

Demonstramos que $A^n$ é a matriz nula. Pela definição de uma matriz nilpotente, encontramos um
inteiro positivo, a saber $k=n$, para o qual $A^k=0$.

Portanto, a matriz $A$ é nilpotente.

\subsection*{Quod Erat Demonstrandum (Q.E.D.)}

\subsection*{Observações Finais}

A beleza desta prova reside na sua abstração. Não precisamos em momento algum conhecer as
entradas da matriz $A$. Apenas com a informação sobre seu espectro de autovalores, e armados com
a ferramenta correta — o Teorema de Cayley-Hamilton —, a estrutura da matriz se revela.

Uma prova alternativa, igualmente instrutiva, poderia ser construída utilizando a Forma Canônica de
Jordan. Como o único autovalor é $0$, a forma de Jordan de $A$ seria uma matriz estritamente triangular
superior (zeros na diagonal principal). É um exercício clássico provar que qualquer matriz
estritamente triangular superior de ordem $n$ elevada à potência $n$ resulta na matriz nula. Como $A$ é
semelhante à sua forma de Jordan $J$ (i.e., $A = PJP^{-1}$), teríamos $A^n = (PJP^{-1})^n =
P J^n P^{-1} = P \cdot 0 \cdot P^{-1} = 0$. Ambas as abordagens levam à mesma conclusão inelutável.

Estudem esta demonstração. Ela é um exemplo primoroso de como os conceitos teóricos que
desenvolvemos nos permitem resolver problemas concretos de maneira concisa e poderosa.